% mr_t4_oos — Out-of-sample summary: recursive Campbell-Thompson R2 by phase.
% Numbers transcribed from tables/mr_t4_oos.pdf (generated by main_results.R).
\small
\begin{tabular}{ll ccc}
  \toprule
  Phase & Specification & In-sample $R^{2}$ (mean) & $R^{2}_{\mathrm{oos}}$ (pooled) & Markets OOS$+$ \\
  \midrule
  I --- local          & $\rxbar \sim \CF$ & $0.254$ & $-0.073$ & $2/11$ \\
  II --- global        & $\rxbar \sim \GCF$ & $0.254$ & $0.048$  & $10/11$ \\
  III --- USD, global  & $\rxbar^{\,\mathrm{USD}} \sim \GCF$ & $0.086$ & $-0.070$ & $2/11$ \\
  III --- USD, FX-adj. & $\rxbar^{\,\mathrm{USD}} \sim \FXGCF$ & $0.143$ & $0.021$ & $6/11$ \\
  \bottomrule
\end{tabular}
\tabnotes{The in-sample $R^{2}$ is the cross-country mean of the single-factor
  fits \eqref{eq:h-local}, \eqref{eq:h-global}, \eqref{eq:h-usd-global} and
  \eqref{eq:h-fxgcf}, i.e.\ variance explained around the full-sample mean.
  The out-of-sample $R^{2}_{\mathrm{oos}}$ is the pooled
  \citet{campbellthompson2008} statistic
  $1-\mathrm{SSE}_{\text{factor}}/\mathrm{SSE}_{\text{mean}}$ of the fully
  recursive factor forecast relative to the recursive prevailing-mean forecast
  \eqref{eq:meth-ctr2}; a positive value means the factor beats the real-time
  historical average. Pooling sums each country's squared forecast errors
  against its own recursive mean before forming the ratio, and ``Markets
  OOS$+$'' counts the markets with a positive $R^{2}_{\mathrm{oos}}$. Both the
  factor construction and the forecasting regression respect time-$t$
  information, so the forecasts are doubly out of sample. The two $R^{2}$
  columns use different benchmarks and are not comparable in level.
  Per-country detail is developed in \Cref{ch:robustness}.}
